\documentclass[
    12pt,
    openright,
    oneside,
    a4paper,
    chapter=TITLE,
    section=TITLE,
    brazil,
    sumario=tradicional,
    hidelinks
]{abntex2}

\usepackage[utf8]{inputenc}
\usepackage[T1]{fontenc}
\usepackage{iftex}
\ifPDFTeX
\usepackage[scaled]{helvet}
\renewcommand\familydefault{\sfdefault}
\else
\usepackage{fontspec}
\setmainfont{Arial}
\setsansfont{Arial}
\fi
\usepackage{microtype}
\usepackage{longtable}
\usepackage{array}
\usepackage{booktabs}
\usepackage{enumitem}

\setlength{\parindent}{1.25cm}
\setlength{\parskip}{0.2cm}
\OnehalfSpacing
\setlist[itemize]{leftmargin=1.2cm}
\newcolumntype{L}[1]{>{\raggedright\arraybackslash}p{#1}}

\titulo{Relatório Técnico de Entregas Trimestrais - HEAL+ / REDI-SUS}
\autor{Equipe HEAL+ (Fatec Ferraz)}
\local{Suzano - SP}
\data{Março de 2026}
\instituicao{%
    Faculdade de Tecnologia de Ferraz de Vasconcelos\\
  Projeto REDI-SUS
}
\tipotrabalho{Relatório técnico}
\preambulo{Documento de alinhamento interno das entregas trimestrais do módulo de diagnóstico do HEAL+ no contexto do REDI-SUS.}

\begin{document}

\imprimircapa
\imprimirfolhaderosto*

\begin{resumo}
Este relatório detalha, por trimestre (T1 a T8), as entregas previstas para o HEAL+ no contexto do REDI-SUS, com foco no módulo de diagnóstico por imagem e na interação com os pacotes correlatos DermaSUS e REDE VIVA. O documento organiza o escopo por pacote de trabalho (PT), define evidências mínimas por entrega e apresenta uma matriz de interação para alinhamento entre equipes. Também propõe checklist de validação para reunião técnica, critérios transversais de aceite e encaminhamentos para consolidação pós-reunião.

\textbf{Palavras-chave}: HEAL+, REDI-SUS, entregas trimestrais, diagnóstico por imagem, interoperabilidade.
\end{resumo}

\pdfbookmark[0]{\contentsname}{toc}
\tableofcontents*
\cleardoublepage

\textual

\section{Objetivo do documento}
Detalhar, por trimestre (T1 a T8), as entregas do HEAL+ no contexto do REDI-SUS, com foco no módulo de diagnóstico por imagem e na interação com os pacotes correlatos (DermaSUS e REDE VIVA). Este documento é uma base de trabalho para reunião de alinhamento e consolidação de evidências de entrega.

\section{Escopo e recorte HEAL+}
\begin{itemize}
    \item Eixo principal: diagnóstico e monitoramento de feridas com IA (captura, pré-processamento, detecção, segmentação, classificação e suporte à decisão).
    \item Pacotes de trabalho com maior impacto no módulo de diagnóstico: PT2, PT3, PT4, PT5, PT6 e PT7.
    \item Interdependências diretas: padronização de dados clínicos, interoperabilidade FHIR e validação em cenários reais.
\end{itemize}

\section{Mapa trimestral de entregas HEAL+ (T1--T8)}
\noindent\textit{Mapa ajustado para leitura em relatório técnico: colunas com melhor distribuição e quebra de linha automática.}

{\footnotesize
\renewcommand{\arraystretch}{1.25}
\setlength{\tabcolsep}{3pt}
\begin{longtable}{L{0.9cm}L{1.1cm}L{3.2cm}L{4.2cm}L{3.1cm}}
\toprule
\textbf{Trim} & \textbf{PT} & \textbf{Entrega (produto)} & \textbf{Detalhamento do que será entregue} & \textbf{Evidência mínima esperada} \\
\midrule
\endhead

T1 & PT2 & Documento de Requisitos Técnicos e Clínicos do REDI-SUS & Especificação de requisitos funcionais e não funcionais para captura de imagem, trilha de diagnóstico, perfis de usuário e segurança/LGPD. & Documento versionado + matriz de requisitos + rastreabilidade inicial. \\

T2 & PT2 & Relatório Técnico da Arquitetura Conceitual e Dicionário de Dados & Definição da arquitetura do módulo de diagnóstico, contratos de dados e padrões terminológicos (CID-10, SNOMED CT, LOINC quando aplicável). & Relatório técnico + modelo de dados + dicionário de campos clínicos. \\

T3 & PT3 & Protocolos de coleta e calibração de imagens clínicas de feridas & Protocolo de captura (distância, iluminação, enquadramento, qualidade), critérios de exclusão e rotinas de controle de qualidade de imagem. & Protocolo operacional + checklist de coleta + formulário de auditoria. \\

T4 & PT3 & Bases de dados clínicos com imagens rotuladas (Heal+) & Base inicial curada com classes clínicas e metadados padronizados para treinamento/validação do módulo de diagnóstico. & Inventário de base + estatísticas de rótulos + relatório de curadoria. \\

T4 & PT3 & Módulos de coleta de dados (REDE VIVA; Heal+; DermaSUS) & Fluxo de ingestão e pré-processamento para dados multimodais (imagem e dados clínicos associados). & Pipeline executável + registros de execução + especificação técnica. \\

T4 & PT4 & Protótipos iniciais de modelos de IA e relatório de arquiteturas & Protótipo de detecção e classificação de lesões com métricas preliminares e comparativo de arquiteturas. & Modelos treinados (versão inicial) + métricas baseline + relatório técnico. \\

T5 & PT3 & Bases de dados clínicos consolidadas (REDE VIVA; Heal+; DermaSUS) & Consolidação da base multicentro com controle de qualidade, deduplicação e padronização de metadados clínicos. & Relatório de consolidação + versão oficial de dataset + log de qualidade. \\

T5 & PT5 & Protótipos de interface diagnóstica e jornadas validadas & Protótipos de interface do processo diagnóstico para profissional de saúde e validação preliminar de fluxo de uso. & Protótipo navegável + roteiro de teste + devolutiva dos usuários. \\

T5 & PT7 & Gateway FHIR operacional & Interoperabilidade para troca estruturada de dados clínicos e resultados de diagnóstico com sistemas SUS. & Endpoints ativos + testes de integração + documentação API. \\

T6 & PT4 & Biblioteca federada de modelos de IA e módulos de apoio à decisão & Versões de modelos com governança de publicação, metadados de treino e estratégia de atualização/federação. & Catálogo de modelos + artefatos versionados + guia de uso. \\

T6 & PT5 & Protótipos funcionais e relatório de usabilidade & Protótipo funcional do fluxo diagnóstico com avaliação de usabilidade e pontos de melhoria priorizados. & Relatório SUS/NPS/escala de usabilidade + backlog priorizado. \\

T6 & PT6 & Relatório de início dos pilotos multicêntricos e plano de avaliação & Definição dos pilotos, critérios de inclusão, indicadores e plano de monitoramento de desempenho clínico-operacional. & Plano de avaliação + cronograma de piloto + termo de governança. \\

T7 & PT4 & Sistema integrado de IA com painéis preditivos e validação & Integração dos modelos em ambiente de uso e painel de indicadores para apoio à decisão clínica. & Sistema integrado + painel operacional + relatório de validação. \\

T7 & PT5 & Materiais de capacitação digital e percepção do usuário & Materiais de treinamento para profissionais e relatório de satisfação/usabilidade/autoeficácia. & Kit de capacitação + relatório analítico de percepção. \\

T7 & PT7 & Plataforma interoperável com SUS digital e rastreabilidade & Evolução da camada de interoperabilidade com rastreabilidade de eventos clínicos e técnicos. & Dashboard de rastreabilidade + documentação arquitetural. \\

T8 & PT6 & Relatório final de validação e transferência tecnológica & Consolidação dos resultados dos pilotos, limites do modelo e proposta de transferência/escala. & Relatório final + plano de transferência + recomendações de escala. \\
\bottomrule
\end{longtable}
}

\section{Recorte para reunião: módulo de diagnóstico (Heal+ x DermaSUS x REDE VIVA)}
\textbf{Objetivo da reunião}: confirmar o que cada equipe realmente já concluiu, o que está parcialmente pronto e o que ainda precisa de ajuste para aderência ao REDI-SUS.

\subsection{Matriz de interação entre pacotes}
{\footnotesize
\renewcommand{\arraystretch}{1.25}
\setlength{\tabcolsep}{3pt}
\begin{longtable}{L{0.9cm}L{2.7cm}L{3.3cm}L{3.3cm}L{2.8cm}}
\toprule
\textbf{Trim} & \textbf{Entrega foco (diagnóstico)} & \textbf{Contribuição esperada DermaSUS} & \textbf{Contribuição esperada REDE VIVA} & \textbf{Ponto de validação na reunião} \\
\midrule
\endhead

T3 & Protocolo de coleta/calibração & Padrão de captura dermatológica e rotulação de lesões de pele. & Procedimentos de campo e contexto epidemiológico para coleta consistente. & Confirmar protocolo único (ou tabela de equivalência entre protocolos). \\

T4 & Módulos de coleta e base rotulada inicial & Dados dermatológicos curados + revisão de qualidade de rótulo. & Dados de campo e metadados clínico-epidemiológicos com padrão mínimo. & Verificar formato único de metadados e processo de QA compartilhado. \\

T5 & Base consolidada e interface diagnóstica & Validação de classes dermatológicas e testes de uso com especialistas. & Casos reais de uso e validação de fluxo em cenários externos. & Fechar critério de aceite conjunto para dataset e UI. \\

T6 & Biblioteca federada + pilotos & Contribuir com modelos/parâmetros e validação cruzada dermatológica. & Contribuir com cenários de piloto e eventos de monitoramento de desempenho. & Confirmar versão de modelo por parceiro e regra de atualização. \\

T7 & Sistema integrado de IA + painéis & Insumos para análise clínica comparativa e confiabilidade do diagnóstico. & Indicadores de campo para vigilância e tomada de decisão. & Definir KPI final por parceiro e painel único de acompanhamento. \\
\bottomrule
\end{longtable}
}

\subsection{Checklist rápido para coleta de status (amanhã)}
\begin{itemize}
    \item O que já está pronto e demonstrável por equipe (arquivo, protótipo, API, base, relatório)?
    \item O que existe, mas ainda não está aderente ao escopo REDI-SUS?
    \item Quais evidências faltam para considerar a entrega concluída no trimestre?
    \item Quais dependências externas impedem fechamento imediato?
    \item Quem valida tecnicamente cada entrega e em qual prazo?
\end{itemize}

\section{Critérios de aceite transversais (sugestão para o grupo)}
\begin{itemize}
    \item \textbf{Rastreabilidade}: cada entrega deve apontar requisito de origem, artefato e responsável.
    \item \textbf{Reprodutibilidade}: processo de geração de resultados deve poder ser repetido por outra equipe.
    \item \textbf{Interoperabilidade}: dados clínicos e resultados de IA devem seguir contrato de dados comum (FHIR quando aplicável).
    \item \textbf{Qualidade de dados}: bases devem ter documentação de cobertura, distribuição de classes e controle de inconsistências.
    \item \textbf{Usabilidade e segurança}: para entregas de interface, incluir evidências de teste de usabilidade e requisitos de privacidade/LGPD.
\end{itemize}

\section{Encaminhamento pós-reunião}
\begin{itemize}
    \item Consolidar status por trimestre: \textit{Concluído}, \textit{Parcial}, \textit{Não iniciado}, \textit{Fora de escopo}.
    \item Atualizar cronograma com riscos, dependências e replanejamento de marcos críticos.
    \item Gerar versão final do relatório trimestral com anexos de evidências técnicas.
\end{itemize}

\vspace{0.5cm}
\noindent\textbf{Observação}: este documento foi estruturado para apoiar o alinhamento interno e deve ser atualizado após a reunião com DermaSUS e REDE VIVA, incorporando o status real de cada artefato.

\postextual

\end{document}
